% Chapter 00: Introduction
% What Domain Kits are, why they exist, who this guide is for

\chapter{Introduction}
\label{ch:introduction}

% ============================================================
\section{What Are Domain Kits?}
\label{sec:what-are-kits}
% ============================================================

The gert v2 runtime is deliberately domain-agnostic. It executes steps, enforces
governance, captures evidence, and writes traces---but it carries no opinion
about \emph{what kind of work} a runbook describes. Whether you are orchestrating
a compliance audit, provisioning infrastructure, migrating databases, or
automating incident response, the gert core executes the same primitive
operations: run CLI commands, invoke tools, pause for manual approval, branch on
conditions, iterate over lists.

This domain neutrality is not an omission---it is the core design principle.
A stable, long-lived runtime cannot accommodate an ever-expanding union of
domain-specific vocabularies without becoming brittle and unmaintainable.

\textbf{Domain Kits} are the mechanism by which domain-specific semantics enter
gert without contaminating the kernel. A Domain Kit is a compile-time semantic
layer that sits \emph{above} the gert core runtime. It provides:

\begin{enumerate}
  \item \textbf{Authoring schemas} --- domain-specific YAML vocabulary (custom
    step types, field structures, validation constraints) that feels natural for
    authors working in a particular domain.
  \item \textbf{Validation rules} --- semantic validators that enforce
    domain-specific invariants (``every approval must have a fallback
    timeout,'' ``cleanup must follow provision'').
  \item \textbf{Lowering (compilation)} --- a pure transformation that compiles
    Kit-specific authoring syntax into core gert primitives before execution.
  \item \textbf{Projections} --- read-only transformations over run traces that
    produce domain-specific reports (audit logs, timeline views, evidence
    bundles).
  \item \textbf{Testing contracts} --- fixture runbooks and expected outputs
    that define what ``correct'' means for this domain.
\end{enumerate}

Domain Kits are \emph{not} runtime extensions. They do not add new capabilities
to the gert engine, register tools, or subscribe to events. A Kit operates
entirely at \textbf{author time} and \textbf{compile time}. By the time the
gert runtime starts executing a runbook, all Kit-specific syntax has been
lowered into core primitives.

% ============================================================
\section{Why Domain Kits Exist}
\label{sec:why-kits}
% ============================================================

Without Domain Kits, domain-specific semantics have two bad places to live:

\paragraph{Option 1: Leak into the core schema.}
Every new domain (compliance, workflow automation, data migration, system
provisioning) adds fields and step types to the core. The schema becomes a
growing union type. The parser grows mode flags. The runtime accumulates
conditionals. The kernel loses its stability guarantee. This is the path to a
bloated, unmaintainable system.

\paragraph{Option 2: Push into extensions.}
Extensions add runtime capabilities (tools, event subscriptions, policy
enforcement). If we also use them for authoring semantics, we conflate two
fundamentally different concerns. An extension that contributes validators,
authoring schemas, and projections is doing work that has nothing to do with
runtime capabilities. Worse, it forces authoring-layer concerns into a runtime
architecture, making them heavier than they need to be.

\paragraph{Domain Kits solve this by providing a clean home for domain
semantics.} They preserve the kernel's stability while allowing unlimited
domain-specific expressiveness. The gert core schema never changes when a new
domain is added---the Kit compiles domain vocabulary down to core primitives.

% ============================================================
\section{Relationship to Gert Core}
\label{sec:kit-and-core}
% ============================================================

Domain Kits \textbf{compile down to core primitives}. The gert core is
unchanged. A Kit is a higher-level authoring syntax built \emph{on top of} the
core runtime, not a modification of it.

The pipeline is:

\begin{verbatim}
Kit YAML → Kit Parser → Kit AST → Kit Compiler → Core YAML → Core Parser → ExecutionPlan
\end{verbatim}

Key properties:

\begin{itemize}
  \item The Kit compiler is a \textbf{pure function}: YAML in, YAML out, no
    side effects.
  \item The lowered output is a valid \texttt{runbook/v2} document that uses
    only core step types (\texttt{cli}, \texttt{tool}, \texttt{manual},
    \texttt{branch}, \texttt{iterate}).
  \item The runtime \textbf{never sees Kit-specific nodes}. Lowering happens
    before the planner produces an ExecutionPlan.
  \item The lowered runbook can be inspected, version-controlled, and executed
    \textbf{without the Kit installed}. This is critical for portability and
    auditability.
\end{itemize}

% ============================================================
\section{The Three Extension Points}
\label{sec:three-extension-points}
% ============================================================

Gert v2 has three distinct extension mechanisms. Each serves a different purpose
and operates at a different phase:

\begin{description}
  \item[Domain Kits (authoring semantics)] contribute domain-specific YAML
    vocabulary, validation rules, lowering logic, and projections. They operate
    at \textbf{author time} and \textbf{compile time}. They are pure
    transformations with no runtime footprint.

  \item[Extensions (runtime capabilities)] contribute tools, event
    subscriptions, schema field extensions (\texttt{x-*} namespaced fields),
    policy rules, and provider integrations. They operate at \textbf{run time}
    as out-of-process children of the gert engine.

  \item[Tools (effectful actions)] perform side-effecting work: binary
    execution, API calls, MCP tool invocations. They operate at \textbf{run
    time}, invoked on demand by the gert engine when a \texttt{tool} step
    executes.
\end{description}

\paragraph{How Kits differ from Extensions and Tools.}

\begin{itemize}
  \item \textbf{Kits are compile-time; Extensions are runtime.} A Kit compiler
    runs once before execution. An Extension runs throughout execution as a
    persistent child process.
  \item \textbf{Kits are pure; Extensions and Tools have side effects.} A Kit
    never invokes binaries, calls APIs, or writes to disk (except for its own
    compiled output). Extensions and Tools perform effectful work.
  \item \textbf{Kits shape vocabulary; Extensions shape capabilities.} A Kit
    determines what authors can write. An Extension determines what the runtime
    can do.
\end{itemize}

A Kit distribution may \emph{ship alongside} an Extension (e.g., a ``compliance
pack'' that includes both a Domain Kit and an Extension with compliance-specific
tools), but they remain architecturally distinct.

% ============================================================
\section{Who Should Read This Guide}
\label{sec:audience}
% ============================================================

This guide is for \textbf{Kit developers}---engineers who are building
domain-specific authoring layers for gert. It is \emph{not} for runbook authors
who are using an existing Kit.

You should read this guide if you:

\begin{itemize}
  \item Want to create a custom step type for your domain (e.g.,
    \texttt{approval-request}, \texttt{deploy-canary}, \texttt{migrate-table}).
  \item Need to enforce domain-specific validation rules that cannot be
    expressed in JSON Schema alone (e.g., ``every provision step must be
    followed by a cleanup step'').
  \item Want to generate domain-specific reports from gert run traces (e.g.,
    audit logs, compliance evidence bundles, timeline views).
  \item Are building a reusable authoring framework for a community (e.g., a
    Kit for Kubernetes deployments, database migrations, or security audits).
\end{itemize}

\paragraph{Prerequisites.}
This guide assumes you are comfortable with:

\begin{itemize}
  \item YAML syntax and structure
  \item JSON Schema (Draft 2020-12)
  \item Writing and invoking command-line tools
  \item Basic compiler concepts (parsing, transformation, lowering)
\end{itemize}

You do \emph{not} need to know the gert v2 runtime internals. This guide stands
alone. However, for deeper understanding of the core execution model, consult
the \emph{Gert v2 Design Document} (particularly Chapter~4, Domain Kit Model).

% ============================================================
\section{How This Guide Is Structured}
\label{sec:structure}
% ============================================================

This guide follows the Kit development lifecycle:

\begin{description}
  \item[Chapter 1: Kit Anatomy] --- directory structure, manifest format,
    versioning contract.
  \item[Chapter 2: Authoring Schemas] --- how to define custom step types as
    YAML overlays on core schemas.
  \item[Chapter 3: Validation Rules] --- writing parse-time and plan-time
    validators.
  \item[Chapter 4: Lowering] --- writing the Kit compiler that transforms your
    authoring syntax into core primitives.
  \item[Chapter 5: Projections] --- building read models over run traces.
  \item[Chapter 6: Testing Contracts] --- acceptance tests for your Kit.
  \item[Chapter 7: Publishing] --- packaging and distributing your Kit.
  \item[Chapter 8: Reference] --- complete field reference for manifest and
    binary interfaces.
\end{description}

Work through the chapters sequentially if you are building a Kit from scratch.
Use Chapter~8 as a quick reference once you are familiar with the model.

% ============================================================
\section{Forward Reference}
\label{sec:forward-ref}
% ============================================================

For details on the gert v2 runtime, execution model, governance layer, evidence
system, and event model, consult the \emph{Gert v2 Design Document}. In
particular:

\begin{itemize}
  \item \textbf{Chapter 4: Domain Kit Model} --- the architectural rationale
    for Kits and their relationship to core primitives.
  \item \textbf{Chapter 5: Runtime Architecture} --- the execution loop, state
    machine, and planner.
  \item \textbf{Chapter 6: Governance and Policy} --- allowlists, redaction,
    approval gates, evidence capture.
  \item \textbf{Chapter 9: Event Model} --- trace events, event schema, and
    event lifecycle.
\end{itemize}

This guide focuses exclusively on \emph{building} Kits, not on the core runtime
that executes them.
